% !TeX program = pdflatex
%
% A standalone, compact specification of the verifier algorithm implemented by
%   fixtures/moonlight-wrap/Halo2Verifier.sol
%   fixtures/moonlight-wrap/Halo2VerifyingKey.sol
%
% This document is deliberately independent of docs/spec/Halo2VerifierSpec.tex.
% That document is a line-by-line, excerpt-quoting reference (~19k lines); this
% one states the algorithm and the complete check catalogue, and is meant to be
% readable end-to-end in one sitting.
%
% Build: pdflatex MoonlightWrapVerifierAlgorithm.tex  (three passes for refs)
% Packages: all ship with texlive-latex-recommended + texlive-science.

\documentclass[11pt,a4paper]{article}

\usepackage[T1]{fontenc}
\usepackage[utf8]{inputenc}
\usepackage[margin=2.4cm]{geometry}
\usepackage{amsmath}
\usepackage{amssymb}
\usepackage{amsthm}
\usepackage{booktabs}
\usepackage{longtable}
\usepackage{array}
\usepackage{enumitem}
\usepackage{xcolor}
\usepackage[hidelinks]{hyperref}

% ---------------------------------------------------------------------------
% Notation macros
% ---------------------------------------------------------------------------
\newcommand{\Fr}{\ensuremath{\mathbb{F}_{r}}}
\newcommand{\Fp}{\ensuremath{\mathbb{F}_{p}}}
\newcommand{\Fptwo}{\ensuremath{\mathbb{F}_{p^{2}}}}
\newcommand{\Gone}{\ensuremath{\mathbb{G}_{1}}}
\newcommand{\Gtwo}{\ensuremath{\mathbb{G}_{2}}}
\newcommand{\GT}{\ensuremath{\mathbb{G}_{T}}}
\newcommand{\Ident}{\ensuremath{\mathcal{O}}}
\newcommand{\pair}{\ensuremath{e}}
\newcommand{\trunc}{\ensuremath{\tau}}
\newcommand{\code}[1]{\texttt{#1}}
\newcommand{\err}[1]{\textcolor{red!60!black}{\texttt{#1}}}
\newcommand{\mptr}[1]{\texttt{#1}}
\newcommand{\opc}[1]{\textsc{#1}}

\theoremstyle{definition}
\newtheorem{remark}{Remark}[section]

\title{The Moonlight-Wrap Halo2/KZG Solidity Verifier\\[2mm]
\large Algorithm, algebra, and complete check catalogue}
\author{Independent extraction from the rendered artifacts}
\date{Fixture render of 2026-08-20 \quad(\code{fixtures/moonlight-wrap/})}

\begin{document}
\maketitle

\begin{abstract}
\noindent
This document specifies what
\code{fixtures/moonlight-wrap/Halo2Verifier.sol} (4{,}130 lines) and
\code{fixtures/moonlight-wrap/Halo2VerifyingKey.sol} (695 lines) compute and
what they reject. It is a \emph{circuit-specialised} Halo2 verifier over
BLS12-381 with a KZG polynomial commitment scheme, targeting the EIP-2537
precompiles, generated by this repository's Rust code generator from one fixed
\code{VerifyingKey}. Section~\ref{sec:prelim} fixes notation and encodings;
Section~\ref{sec:artifacts} describes the two-contract split and how each half
pins the other; Sections~\ref{sec:deploy}--\ref{sec:pairing} give the verifier
algorithm in order of execution together with the algebra of each phase; and
Section~\ref{sec:catalogue} is a complete catalogue of every runtime check and
the typed error it raises. Section~\ref{sec:notchecked} states, explicitly,
what the contract does \emph{not} check and where that obligation lives
instead.

All structural constants quoted here (proof layout, identity count, point
sets, MSM width) were read out of the rendered artifacts, not out of the
generator.
\end{abstract}

\tableofcontents

% ===========================================================================
\section{Preliminaries and notation}
\label{sec:prelim}
% ===========================================================================

\subsection{Fields and groups}

Let $E/\Fp: y^{2}=x^{3}+4$ be BLS12-381 with
\[
p = \mathtt{0x1a0111ea\ldots abaaab}\ (381\ \text{bits}),
\qquad
r = \mathtt{0x73eda753299d7d483339d80809a1d80553bda402fffe5bfeffffffff00000001}\ (255\ \text{bits}).
\]
$\Gone \subset E(\Fp)$ and $\Gtwo \subset E'(\Fptwo)$ are the order-$r$
subgroups, and $\pair:\Gone\times\Gtwo\to\GT$ is the optimal Ate pairing as
exposed by the EIP-2537 precompile \code{0x0f}. Every scalar the verifier
manipulates lives in $\Fr$; coordinates live in $\Fp$. The main verifier body binds
$r$ once (\code{let r := FR\_MODULUS}; Yul helper functions rebind the same
constant locally) and every arithmetic line is an
\code{addmod}/\code{mulmod} modulo that $r$.

Write $\Ident$ for the point at infinity of $\Gone$.

\subsection{Trusted setup}

The verifying key carries three points, all taken on trust:
\[
G \in \Gone \quad(\text{the generator}),\qquad
[1]_2 \in \Gtwo,\qquad
[-s]_2 \in \Gtwo,
\]
where $s$ is the KZG toxic waste. Only $[1]_2$ and $[-s]_2$ enter the final
pairing; $[s]_2$ never appears explicitly.

\subsection{Encodings}

\paragraph{EIP-2537 padded coordinates.}
An $\Fp$ element is transmitted as two 32-byte words
\[
c \;=\; c_{\mathrm{hi}}\cdot 2^{256} + c_{\mathrm{lo}},
\qquad
c_{\mathrm{hi}} < 2^{128},
\]
i.e.\ 16 zero pad bytes followed by the 48 big-endian coordinate bytes. A
$\Gone$ point is therefore $4$ words / $128$ bytes
$(x_{\mathrm{hi}},x_{\mathrm{lo}},y_{\mathrm{hi}},y_{\mathrm{lo}})$ and a
$\Gtwo$ point is $8$ words / $256$ bytes. The identity $\Ident$ is encoded as
four zero words.

\paragraph{Truncation.}
Define the 128-bit truncation
\[
\trunc(a) \;=\; a \bmod 2^{128}.
\]
The verifier applies $\trunc$ to $x_3$ immediately after squeezing it, and to
the \emph{stored copies} of the powers $x_1^{i}$ and $x_4^{i}$ while keeping
the internal power accumulators at full precision. This mirrors the
truncated-challenge convention of \code{midnight-proofs}'
\code{poly/kzg/mod.rs}.

\paragraph{Scalars.}
Public instances, proof evaluations and \code{q\_eval}s are transmitted as
canonical big-endian $\Fr$ words; anything ${\ge}\,r$ is rejected. The proof's
$\Gone$ elements are transmitted in the padded \emph{uncompressed} form above:
an off-chain shim repacks \code{midnight-proofs}' native compressed stream
before calling \code{verifyProof}.

\subsection{Circuit parameters of this render}

Every number below is baked into the artifact. Six of the eleven VK header
words are cross-checked against generated constants at run time
(Section~\ref{sec:vkload}); everything else is covered by the VK codehash pin
alone.

\begin{center}
\begin{tabular}{@{}llr@{}}
\toprule
Symbol & Meaning & Value \\
\midrule
$k$                & $\log_2$ domain size            & $20$ \\
$n = 2^{k}$        & domain size                     & $1{,}048{,}576$ \\
$\omega$           & primitive $n$-th root of unity  & VK word 4 \\
$\ell$             & $|{\text{last rotation}}|$, $\omega^{-\ell}$ in VK word 6 & $10$ \\
$N_{\mathrm{inst}}$& public instances                & $19$ \\
$N$                & identities in the $y$-batch     & $49$ \\
$N_{\mathrm{sel}}$ & simple selectors (linearised)   & $10$ \\
--                 & permutation columns / sets / chunk & $18 / 6 / 3$ \\
--                 & lookup arguments                & $2$ \\
--                 & quotient limbs                  & $4$ \\
--                 & distinct rotations              & $4$ \\
--                 & KZG point sets                  & $5$ \\
--                 & terms in the fused final MSM    & $78$ \\
\bottomrule
\end{tabular}
\end{center}

The accumulator (IVC) parameters are
$\mathrm{acc\_offset}=11$, $\mathrm{num\_acc\_limbs}=7$,
$\mathrm{num\_acc\_limb\_bits}=56$.

\subsection{Proof layout}

\code{verifyProof(bytes proof, uint256[] instances)} accepts exactly
$\code{calldatasize} = \mathtt{0x2144}$ bytes: a 4-byte selector, two ABI
heads, two array-length words, a $\mathtt{0x1e60}=7776$-byte \code{proof}
payload, and $19$ instance words. The $7776$ proof bytes are $34$ padded $\Gone$ points ($4352$ bytes)
and $107$ scalars ($3424$ bytes), laid out in the order the transcript
consumes them:

\begin{center}
\begin{tabular}{@{}rlr@{}}
\toprule
\# & Item & Bytes \\
\midrule
15 & advice commitments (single user phase) & $1920$ \\
 2 & lookup multiplicity commitments $M_c$  & $256$ \\
 6 & permutation product commitments $Z_s$  & $768$ \\
 4 & lookup helper/accumulator pairs, interleaved per argument ($H_0, Z^{L}_0, H_1, Z^{L}_1$) & $512$ \\
 1 & trashcan commitment                    & $128$ \\
 4 & quotient limb commitments $Q_0..Q_3$   & $512$ \\
102 & polynomial evaluations at rotations of $x$ & $3264$ \\
 1 & batched-opening commitment $F$          & $128$ \\
 5 & per-point-set evaluations $q_0..q_4$ at $x_3$ & $160$ \\
 1 & KZG opening proof $\pi$                 & $128$ \\
\midrule
   & \textbf{total}                          & $\mathbf{7776}$ \\
\bottomrule
\end{tabular}
\end{center}

\noindent
The parser walks these with a single cursor and, after the last item, asserts
that the cursor landed exactly on the \code{instances} length word --- a
whole-layout consistency check that is redundant with the length constants but
fails closed if either drifts.

% ===========================================================================
\section{The two artifacts and how they pin each other}
\label{sec:artifacts}
% ===========================================================================

\subsection{\code{Halo2VerifyingKey}: data as bytecode}

The verifying key is a contract whose \emph{runtime bytecode} is
\[
\mathtt{0xfe} \;\|\; \underbrace{\text{payload}}_{17{,}024\ \text{bytes} \;=\; 532\ \text{words}},
\qquad
L_{\mathrm{runtime}} = 17{,}025 .
\]
Byte $0$ is an unconditional \opc{invalid} opcode, so a direct call to the
address cannot execute the payload as code. The constructor writes the blob at
memory $\mathtt{0x80}$ and returns it; there are no functions, no storage, and
no logic. The payload word layout is

\begin{center}
\begin{tabular}{@{}llp{7.2cm}@{}}
\toprule
Words & Contents & Note \\
\midrule
$0$        & $\mathrm{vk\_digest}$   & transcript representation of the constraint system \\
$1$--$2$   & $N_{\mathrm{inst}}$, $k$ & cross-checked at run time \\
$3$--$6$   & $n^{-1}$, $\omega$, $\omega^{-1}$, $\omega^{-\ell}$ & \emph{not} cross-checked \\
$7$--$10$  & accumulator metadata     & cross-checked at run time \\
$11$--$14$ & $G$                      & padded $\Gone$ \\
$15$--$22$ & $[1]_2$                  & padded $\Gtwo$ \\
$23$--$30$ & $[-s]_2$                 & padded $\Gtwo$ \\
$31$--$208$ & quotient constant table & $178$ $\Fr$ words \\
$209$--$351$ & packed quotient bytecode & $4{,}559$ bytes, $101$ instructions \\
$352$--$\ldots$ & fixed commitments    & $4$ words each \\
$\ldots$    & permutation commitments  & $4$ words each, $18$ of them \\
\bottomrule
\end{tabular}
\end{center}

\noindent
In total the payload carries $46$ $\Gone$ points and $2$ $\Gtwo$ points.

\subsection{The pin}

The verifier stores the VK address as an \code{immutable} and pins the runtime
by \emph{both} length and \opc{extcodehash}:
\[
\code{extcodesize}(vk) = 17{,}025
\;\wedge\;
\code{extcodehash}(vk) = \mathtt{0xe68d\ldots cf52}.
\]
This is checked once in the constructor and \emph{again on every proof},
before any of the proof stream is parsed (only the two ABI offset heads are
read earlier). The payload is then
\code{extcodecopy}'d from byte $1$ into memory at $\mptr{VK\_MPTR}$
$=\mathtt{0x3680}$, so all downstream code addresses it as if it had been
embedded.

Because the quotient program lives inside this pinned blob, no proof calldata
can influence the interpreter's control flow: the bytecode is as immutable as
the contract code itself.

\subsection{Build identity}

The verifier exposes
\code{BUILD\_ID = keccak256(domain $\|$ feature profile $\|$ vk\_digest $\|$
VK codehash $\|$ SRS fingerprint $\|$ provenance)}, so a deployment record can
be recomputed by a third party.

% ===========================================================================
\section{Deployment-time checks}
\label{sec:deploy}
% ===========================================================================

The constructor runs \code{require\_eip2537\_precompiles()} \emph{before}
storing the VK address. Its purpose is to make a chain that cannot execute
this verifier fail at deployment rather than on the first proof. Every probe
but one forwards the \emph{exact} gas bound the runtime forwards, so a
repriced precompile schedule is caught here too; the exception is the negative
\code{G1MSM} probe (D6), which forwards a deliberately generous fixed
$200{,}000$ because a rejecting precompile consumes everything forwarded.

\begin{enumerate}[label=\textbf{D\arabic*.},leftmargin=2.2cm]
\item \textbf{Memory layout.} $\code{mload}(\mathtt{0x40}) \le \mathtt{0x1000}$.
      solc's stack-spill reservation must stay below the generated arena base.
      This is the one probe with a typed error, \err{MemoryLayoutViolated()},
      because it reports a \emph{build} fault, not a chain fault.
\item \textbf{\opc{mcopy}.} A store/copy/read round-trip, so a pre-Cancun fork
      fails here.
\item \textbf{\code{modexp} (\code{0x05}).} Known answer using the runtime's own
      operation: $2^{\,r-2} \bmod r$, accepted iff
      $2\cdot\mathrm{result} \equiv 1 \pmod r$. Checking multiplicatively
      rather than against a rendered constant rejects a stub that echoes zeros.
\item \textbf{\code{G1ADD} (\code{0x0b}).} $\Ident+\Ident=\Ident$ (shape), then
      the known answer $G+G=[2]G$ compared limb-for-limb.
\item \textbf{\code{G1MSM} (\code{0x0c}), positive.} $[2]\cdot G = [2]G$.
\item \textbf{\code{G1MSM}, negative.} A point that is on $E(\Fp)$ but
      \emph{not} in the order-$r$ subgroup must be \emph{rejected}. This is the
      only inverted guard in the function, and it is the load-bearing one:
      the runtime performs no curve or subgroup check of its own
      (Section~\ref{sec:notchecked}), delegating both to \code{0x0c}.
\item \textbf{\code{PAIRING\_CHECK} (\code{0x0f}), positive.}
      $\pair(G,H)\cdot \pair(-G,H)=1$, compared strictly against the word $1$.
\item \textbf{\code{PAIRING\_CHECK}, negative.} The same frame with the second
      $\Gone$ sign flipped must return $0$. A chain whose \code{0x0f} always
      returns $1$ would accept every proof, so this probe is as load-bearing
      as D6.
\item \textbf{Full-width \code{G1MSM}.} A $78$-pair all-identity MSM at the
      runtime's exact gas bound, so the largest call the verifier ever makes is
      priced at deployment; accepted only with a $128$-byte return encoding
      $\Ident$.
\item \textbf{All-identity \code{PAIRING\_CHECK}.} One more two-pair pairing
      over all-identity inputs at the runtime's two-pair gas bound, accepted
      only with a $32$-byte return equal to $1$ --- the exact frame shape of
      the runtime's final gate.
\item \textbf{VK pin.} $\code{code.length}$ and $\code{codehash}$ as above.
\end{enumerate}

\begin{remark}
Probes D2--D10 revert bare (\code{revert(0,0)}), by design: they report a
chain-capability failure and the deployment transaction identifies itself. The
cost is that a failed deployment does not say \emph{which} capability is
missing. Every probe also validates \code{returndatasize} before comparing
values, so a stub that returns short data cannot pass.
\end{remark}

% ===========================================================================
\section{\texorpdfstring{$\code{verifyProof}$}{verifyProof}: shape of the whole computation}
\label{sec:overview}
% ===========================================================================

The function is \emph{success-or-revert}: it returns \code{true} or it reverts
with one of eight typed selectors. It never returns \code{false}, so
\code{if (!v.verifyProof(...))} is dead code at every call site.

\begin{center}
\begin{tabular}{@{}rlp{8.4cm}@{}}
\toprule
Phase & Name & Output \\
\midrule
0 & ABI \& memory guards       & --- \\
1 & VK load and envelope       & VK payload at $\mptr{VK\_MPTR}$ \\
2 & Public accumulator         & $\mathrm{acc}_L,\ \mathrm{acc}_R \in \Gone$ \\
3 & Transcript / Fiat--Shamir  & $\theta,\beta,\gamma,\tau,y,x,x_1,x_2,x_3,x_4$; commitments and the $102$ evaluations in memory ($q_0..q_4$ stay in calldata behind a saved cursor) \\
4 & Domain evaluation          & $x^{n}$, $L_0(x)$, $L_{\mathrm{last}}(x)$, $L_{\mathrm{blind}}(x)$, $\mathrm{inst}(x)$ \\
5 & Quotient VM                & $-\nu_y(x)$ and $B_0..B_9$ \\
6 & Linearisation scalars      & $x^{n-1}$, $1-x^{n}$ \\
7 & KZG multi-open             & $v$, $\mathrm{final\_com}$ \\
8 & Pairing (+ accumulator batching) & accept / revert \\
\bottomrule
\end{tabular}
\end{center}

\noindent
Phases run in this order for a reason stated in the source: the accumulator is
validated \emph{before} the transcript so that malformed public inputs cannot
influence challenge derivation, and the final pairing is the only accepting
gate.

% ---------------------------------------------------------------------------
\subsection{Phase 0 --- ABI and memory guards}
% ---------------------------------------------------------------------------

Two guards run before any generated memory is touched.

\begin{align}
\text{(head shape)}\quad
&\code{calldataload}(\mathtt{0x04}) = \mathtt{0x40}
\;\wedge\;
\code{calldataload}(\mathtt{0x24}) = \mathtt{0x1ec0},
\\[1mm]
\text{(spill separation)}\quad
&\code{mload}(\mathtt{0x40}) \le \mptr{TRANSCRIPT\_MPTR} = \mathtt{0x1000}.
\end{align}

The first fixes both ABI offset heads and raises \err{BadCalldataShape()}. The
second raises \err{MemoryLayoutViolated()}: the assembly body writes absolute
addresses upward from $\mathtt{0x1000}$ and never consults the free-memory
pointer, so if a recompile ever pushed solc's spill reservation into that
region no input could verify. Both write their selector inline, because the
\code{fail} helper is not yet in scope --- and they live in different blocks:
the head guard sits in a small non-terminal assembly block of its own, before
the VK immutable is even read, while the spill guard opens the terminal
\code{memory-safe} block that runs the rest of the verifier.

\begin{remark}[The \code{memory-safe} annotation]
The main block is annotated \code{assembly ("memory-safe")} although it writes
memory it never allocated through the free-memory pointer. Three properties
make that sound \emph{here} and all three must hold together: the block is
terminal (no Solidity runs after it), the pragma is pinned so codegen cannot
shift underneath it, and guard~(2) fails closed if the spill region ever
reaches the arena. Lifting this body into a non-terminal context, or floating
the pragma, invalidates the annotation.
\end{remark}

% ---------------------------------------------------------------------------
\subsection{Phase 1 --- Verifying-key load and ABI envelope}
\label{sec:vkload}
% ---------------------------------------------------------------------------

\begin{enumerate}[label=\arabic*.,leftmargin=1cm]
\item Re-pin the VK: $\code{extcodesize}(vk)=17{,}025$ and $\code{extcodehash}(vk)=\mathtt{0xe68d\ldots cf52}$, else
      \err{VkMismatch()}. Re-checking on every proof, not only in the
      constructor, hardens against forks and same-transaction edge cases.
\item \code{extcodecopy}$(vk,\ \mptr{VK\_MPTR},\ 1,\ 17{,}024)$.
\item Cross-check six header words against the generated constants:
      \[
      N_{\mathrm{inst}}=19,\quad k=20,\quad \mathrm{has\_acc}=1,\quad
      \mathrm{acc\_offset}=11,\quad \mathrm{limbs}=7,\quad \mathrm{bits}=56,
      \]
      else \err{VkMismatch()}.
\item Check the ABI envelope:
      \[
      \code{proof.length}=\mathtt{0x1e60},\quad
      \code{instances.length}=19,\quad
      \code{calldatasize}=\mathtt{0x2144},
      \]
      else \err{BadCalldataShape()}.
\end{enumerate}

The exact-\code{calldatasize} equality is what makes every later
\code{calldataload} safe without an individual bound: there is no trailing data
and no zero-extension. It is also why ERC-2771 forwarders and other
calldata-appending relayers cannot call this contract directly.

\begin{remark}
Step 3 is a build-consistency canary, not an independent trust anchor: it
covers $6$ of the $11$ header words. In particular $n^{-1}$, $\omega$,
$\omega^{-1}$ and $\omega^{-\ell}$ --- on which the entire Lagrange and
rotation-point construction rests --- are never algebraically cross-checked at
run time. Their correctness rests solely on the codehash pin.
\end{remark}

% ---------------------------------------------------------------------------
\subsection{Phase 2 --- The public IVC accumulator}
\label{sec:acc}
% ---------------------------------------------------------------------------

Instances $11..18$ (the trailing eight words) encode two $\Gone$ points,
$\mathrm{acc}_L$ and $\mathrm{acc}_R$, carried forward by the recursive KZG
accumulator. Each point is four words: $x$ as two words, then $y$ as two.

\paragraph{The limb codec.}
The circuit exposes an $\Fp$ coordinate through
\code{AssignedField::as\_public\_input} as $n_\lambda = 7$ radix-$2^{56}$
limbs of $(c-1)$, packed $\lambda = 4$ limbs per native $\Fr$ word. With
$\mathrm{coord\_words} = \lceil n_\lambda/\lambda \rceil = 2$, one coordinate
occupies words $w_0, w_1$ and

\begin{equation}
\mathrm{limb}_i \;=\;
\Big\lfloor \frac{\,w_{\lfloor i/\lambda\rfloor} - \delta\cdot[\,i<\lambda\,]\,}{2^{\,56\,(i \bmod \lambda)}} \Big\rfloor
\bmod 2^{56},
\qquad i = 0,\dots,6,
\end{equation}

where $\delta \in \{0, 2^{56}\}$ is the identity-flag adjustment, applied only
to $w_0$. The decoded value is reassembled directly into the EIP-2537 split:

\begin{align}
V_{\mathrm{lo}} &= \Big(\textstyle\sum_{56i < 256} \mathrm{limb}_i\, 2^{56 i}\Big) \bmod 2^{256}, \\
V_{\mathrm{hi}} &= \sum_{\substack{56i<256 \\ 56i+56>256}} \Big\lfloor \frac{\mathrm{limb}_i}{2^{\,256-56i}} \Big\rfloor
             \;+\; \sum_{56i \ge 256} \mathrm{limb}_i\, 2^{\,56i-256}.
\end{align}

For $(56,7)$ the straddling sum ranges over $\{4\}$ alone and the high sum over
$\{5,6\}$.

\paragraph{The shift-by-one convention.}
The codec represents $0$ as $p-1$, so after reassembly
\begin{equation}
c \;=\;
\begin{cases}
0 & \text{if } V = p-1, \\
V + 1 & \text{otherwise.}
\end{cases}
\end{equation}

\paragraph{The identity.}
$\Ident$ has exactly \emph{one} accepted public-input encoding: all four packed
words equal the $p-1$ sentinels, with the identity flag $2^{56}$ added to the
first $x$ word only. It decodes to the four-zero-word EIP-2537 slot. Any other
encoding that would decode to $(0,0)$ is rejected, because EIP-2537 reserves
affine $(0,0)$ for infinity.

\paragraph{Checks.} (a)--(c) run for each of the four coordinates; (d) runs
once per point:
\begin{enumerate}[label=(\alph*),leftmargin=1cm]
\item \textbf{Packing canonicality.} Unused high bits must be zero:
      $w_0 < 2^{224}$ and $w_1 < 2^{168}$. Without this a coordinate would have
      many calldata encodings. As a side effect both bounds sit far below $r$,
      so accumulator instance words are canonical scalars by construction.
\item \textbf{Field range.} $V \le p-1$, tested in split form as
      $V_{\mathrm{hi}} < p_{\mathrm{hi}} \vee (V_{\mathrm{hi}} = p_{\mathrm{hi}} \wedge V_{\mathrm{lo}} \le (p-1)_{\mathrm{lo}})$.
\item \textbf{Padding.} $c_{\mathrm{hi}} < 2^{128}$.
\item \textbf{No unflagged infinity.} $(x,y) \neq (0,0)$ unless the canonical
      identity encoding was used.
\end{enumerate}
All four raise \err{BadPointEncoding()}.

\paragraph{Curve and subgroup membership.}
Are \emph{not} tested here. Instead both decoded points are routed through a
one-pair \code{G1MSM} with scalar $1$:
\[
\mathrm{acc}_L \leftarrow [1]\cdot \mathrm{acc}_L,
\qquad
\mathrm{acc}_R \leftarrow [1]\cdot \mathrm{acc}_R,
\]
whose only purpose is that EIP-2537 rejects off-curve and out-of-subgroup
inputs. The multiplication is performed even for the identity and even for
scalar $1$, precisely so that no point can slip past validation by being
multiplied away. (The two sides are staged asymmetrically: $\mathrm{acc}_L$
as a fixed one-pair call, $\mathrm{acc}_R$ as a variable-length pair list that
supports fixed-base scalar tails and collapses to the same single pair with
scalar $1$ in this render.)

A failing \code{staticcall} is reported as \err{PrecompileFailed()} and a
failing decode as \err{BadPointEncoding()}; the split is correct because each
MSM sits inside an \code{if out} guard, so a decode failure short-circuits
before the precompile flag can be written.

% ---------------------------------------------------------------------------
\subsection{Phase 3 --- The Keccak transcript and the challenge schedule}
\label{sec:transcript}
% ---------------------------------------------------------------------------

The transcript is a streaming byte buffer $T$ based at $\mathtt{0x1000}$ with
two operations:

\begin{align}
\mathrm{common}(b) &: \quad T \mathrel{\|}= b, \\
\mathrm{squeeze}() &: \quad h \leftarrow \mathrm{keccak256}(T); \;\;
T \leftarrow h; \;\; \text{return } h \bmod r .
\end{align}

That is, one Keccak finalisation \emph{reseeds} the buffer with its own digest
and the challenge is the digest reduced modulo $r$. A $\Gone$ commitment is
absorbed as its exact 128 calldata bytes --- matching
\code{Hashable<Keccak256> for G1Projective::to\_input} --- and a scalar as its
canonical 32-byte big-endian word.

The absorption/squeeze order is fixed and mirrors
\code{midfall/proofs/src/plonk/verifier.rs}:

\begin{center}
\begin{tabular}{@{}rlp{7.4cm}@{}}
\toprule
\# & Absorbed / squeezed & Note \\
\midrule
1 & $\mathrm{vk\_digest}$ (32 B)             & commits the constraint system first \\
2 & $128$ zero bytes                          & $\mathrm{committed\_pi}=\Ident$ \\
3 & $19$ (32 B)                               & instance count \\
4 & $\mathrm{inst}_0 \ldots \mathrm{inst}_{18}$ & each checked $< r$ \\
5 & $15$ advice commitments                   & \\
6 & \textbf{squeeze} $\theta$                 & lookup input compression \\
7 & $M_0, M_1$                                & lookup multiplicities \\
8 & \textbf{squeeze} $\beta$, \textbf{squeeze} $\gamma$ & permutation/lookup randomisers \\
9 & $Z_0 \ldots Z_5$                          & permutation products \\
10 & $H_0, Z^{L}_0, H_1, Z^{L}_1$             & per-lookup helper then accumulator \\
11 & \textbf{squeeze} $\tau$                  & trash challenge, squeezed unconditionally \\
12 & trashcan commitment                       & absorbed \emph{after} $\tau$ \\
13 & \textbf{squeeze} $y$                     & identity batching \\
14 & $Q_0 \ldots Q_3$                         & quotient limbs, read only after $y$ \\
15 & \textbf{squeeze} $x$                     & main evaluation point \\
16 & $e_0 \ldots e_{101}$                     & each checked $< r$, spilled to memory \\
17 & \textbf{squeeze} $x_1$, \textbf{squeeze} $x_2$ & multi-open batching \\
18 & $F$                                       & batched-opening commitment \\
19 & \textbf{squeeze} $x_3$, then $x_3 \leftarrow \trunc(x_3)$ & PCS evaluation point \\
20 & $q_0 \ldots q_4$                          & each checked $< r$ \\
21 & \textbf{squeeze} $x_4$                    & final PCS batching \\
22 & $\pi$                                     & opening proof \\
\bottomrule
\end{tabular}
\end{center}

Two ordering facts are load-bearing. $\tau$ is squeezed \emph{before} the
trashcan commitment is absorbed, which is what makes the eight-row trash batch
sound; and $y$ is squeezed before $Q_0..Q_3$ are read, so the quotient
commitments cannot be chosen with knowledge of the batching challenge.

\paragraph{Point encoding checks.}
Every absorbed $\Gone$ point is checked, and only checked, for encoding:
\begin{equation}
x_{\mathrm{hi}} < 2^{128} \;\wedge\; y_{\mathrm{hi}} < 2^{128}
\;\wedge\;
x < p \;\wedge\; y < p,
\end{equation}
raising \err{BadPointEncoding()}. Normalising instead of rejecting would let
several calldata encodings share one transcript. No curve or subgroup check
happens here; see Section~\ref{sec:notchecked}.

\paragraph{Scalar canonicity, and where things land.}
Every absorbed scalar is range-checked against $r$
(\err{NonCanonicalScalar()}), but the timing differs: the $19$ instance words
fold into a deferred success flag checked once after their loop, while the
$102$ evaluations and the five $q_s$ revert immediately per word. Each
absorbed commitment is also \code{calldatacopy}'d into a planned per-category
memory region for the later phases, and the evaluations are spilled to a
reversed frame; the five $q_s$ alone are \emph{not} copied --- the multi-open
reads them from calldata through a cursor saved at absorption time.

\paragraph{Parser closure.}
After $\pi$, the cursor must equal the \code{instances} length word's offset,
else \err{BadCalldataShape()}.

\begin{remark}
$\pi$'s absorption is a no-op for soundness: no challenge is squeezed after it,
so only its \emph{encoding} check is load-bearing. Similarly, challenges are
sampled as $h \bmod r$ from a $256$-bit digest, which is measurably biased ---
$r \approx 2^{254.86}$, so the reduction is not uniform. The bias is far below
any practical soundness threshold but it is a deviation from rejection
sampling.
\end{remark}

% ---------------------------------------------------------------------------
\subsection{Phase 4 --- Domain evaluation}
\label{sec:lagrange}
% ---------------------------------------------------------------------------

All of this phase is pure $\Fr$ arithmetic on the squeezed $x$.

\paragraph{Power of the domain.}
$x^{n}$ is computed by $k=20$ squarings, then $x^{n}-1$.

\paragraph{Lagrange values.}
The verifier needs $29$ Lagrange basis values --- $10$ at negative rotations
and $19$ at the instance rows. Slot $j$ of the denominator scratch
corresponds to the domain point $\omega^{\,j-\ell}$ with $\ell = 10$:
\[
\mathrm{den}_j \;=\; x - \omega^{\,j-\ell},
\qquad j = 0,\dots,28,
\qquad
\mathrm{den}_{29} \;=\; x^{n}-1 .
\]
All $30$ words are inverted in \emph{one} \code{modexp} call by Montgomery's
trick (forward prefix products, one Fermat inversion $T^{\,r-2}$, backward
distribution), and then
\begin{equation}
L_{j-\ell}(x)
\;=\;
\frac{(x^{n}-1)}{n}\cdot\frac{\omega^{\,j-\ell}}{\,x-\omega^{\,j-\ell}\,},
\qquad j = 0,\dots,28 .
\end{equation}

From this run the named values are distilled:
\begin{align}
L_{\mathrm{last}}(x) &= L_{-\ell}(x) &&\text{(slot } 0), \\
L_{\mathrm{blind}}(x) &= \sum_{j=1}^{9} L_{j-\ell}(x) = \sum_{i=-9}^{-1} L_i(x) &&\text{(slots } 1..9), \\
L_0(x) &= \text{slot } 10, \\
\mathrm{inst}(x) &= \sum_{j=0}^{18} L_j(x)\cdot \mathrm{inst}_j &&\text{(slots } 10..28).
\end{align}

\begin{remark}[The slot map is implicit]
Nothing in the artifact states that slot $0$ is $L_{\mathrm{last}}$, slots
$1..9$ are the blinding rows, slot $10$ is $L_0$, and slots $10..28$ are the
instance coefficients. It follows only from $\omega^{-\ell}$ being the seed of
the denominator walk, plus three separately generated offsets
($\mathtt{0x140}$, $\mathtt{0x260}$, $\mathtt{0x3a0}$) that must agree for the
loop to stay inside its scratch. They do agree, exactly and with no slack: the
last instance consumes the last denominator slot.
\end{remark}

\paragraph{Failure taxonomy.}
The batched inversion returns two flags. A failed \code{modexp} (chain fault)
raises \err{PrecompileFailed()}; a zero or non-canonical denominator raises
\err{ProofRejected()}, since the only way to reach it is a squeezed $x$ that
coincides with a domain point --- probability ${\approx}\,n/r$.

\begin{remark}
The named slots are written \emph{before} the success flag is examined, so on a
rejection they briefly hold garbage; the code is safe only because the check
that follows reverts unconditionally before any consumer runs. Two of the six
named slots ($x^{n}$ and $(x^{n}-1)^{-1}$) have no reader anywhere in the file:
$x^{n}$ is recomputed from scratch in Phase 6.
\end{remark}

% ---------------------------------------------------------------------------
\subsection{Phase 5 --- The quotient VM and the batched numerator}
\label{sec:qvm}
% ---------------------------------------------------------------------------

\subsubsection{What is being reconstructed}

The proof does \emph{not} carry a scalar $h(x)$ for the verifier to trust.
Instead the verifier rebuilds the $y$-batched constraint numerator from the
alleged evaluations, and rebuilds the commitment side from the quotient limb
commitments. Let $\mathrm{id}_0,\dots,\mathrm{id}_{N-1}$ be the identities in
the canonical order (gate, permutation, lookup, trash), $N = 49$. Partition
them into
\begin{itemize}[leftmargin=1cm,nosep]
\item \emph{linearised} identities, each of the form
      $\mathrm{id}_i = s_{j(i)}\cdot \widehat{\mathrm{id}}_i$ for a simple
      selector polynomial $s_j$, $j \in \{0,\dots,9\}$, and
\item \emph{fully evaluated} identities, everything else.
\end{itemize}
Define
\begin{align}
\nu_y(x) &\;=\; \sum_{i \text{ fully evaluated}} y^{\,N-1-i}\,\mathrm{id}_i(x),
\label{eq:nu}\\
B_j &\;=\; \sum_{i \in I_j} y^{\,N-1-i}\,\widehat{\mathrm{id}}_i(x),
\qquad j = 0,\dots,9,
\label{eq:bucket}
\end{align}
where $I_j$ is the set of identity positions owned by selector $j$. The
linearised commitment is
\begin{equation}
\mathcal{L}
\;=\;
\sum_{j=0}^{9} B_j \cdot S_j
\;+\;
(1-x^{n})\sum_{\lambda=0}^{3} \big(x^{\,n-1}\big)^{\lambda}\, Q_\lambda,
\label{eq:lincom}
\end{equation}
with $S_j$ the (fixed, VK-resident) selector commitments and $Q_\lambda$ the
quotient limb commitments from the proof, and the claim the PCS must verify is
\begin{equation}
\boxed{\;\mathcal{L}\ \text{opens at}\ x\ \text{to}\ -\nu_y(x).\;}
\label{eq:linclaim}
\end{equation}

Equation~\eqref{eq:linclaim} is exactly the Halo2 quotient identity
$\sum_i y^{N-1-i}\mathrm{id}_i(x) = (x^{n}-1)\,h(x)$ rearranged so that the
untrusted $h(x)$ never appears as a scalar: the $(1-x^n)$ factor supplies the
negated vanishing polynomial and the limb sum supplies $h$'s commitment.

\subsubsection{How $\nu_y(x)$ and the buckets are computed}

Two accumulators are maintained: $A$ (initially $0$), which ends as
$\nu_y(x)$, and $B_0..B_9$ (initially $0$).

\begin{itemize}[leftmargin=1cm]
\item A \emph{fully evaluated} identity contributes
      $A \leftarrow A\cdot y$ then $A \leftarrow A + \mathrm{id}_i(x)$.
\item A \emph{linearised} identity contributes
      $A \leftarrow A\cdot y$ (no addition --- it holds the $y$-position open)
      and $B_j \leftarrow B_j\cdot y^{g} + \widehat{\mathrm{id}}_i(x)$, where
      the gap $g$ is the distance since bucket $j$'s previous contribution
      ($g = 0$ on a bucket's first contribution, skipping the multiply).
\item After the stream, each bucket is aligned to the global end by a
      code-generated tail exponent $t_j$:
      $B_j \leftarrow B_j \cdot y^{\,t_j}$, with
      $t_j = N-1-i_{\mathrm{last}}(j)$.
\end{itemize}
The tail exponents in this render are
$(t_0,\dots,t_9) = (48,47,44,41,38,35,32,29,26,20)$. Since bucket $0$'s only
contribution is at global position $0$, $t_0 = 48$ pins $N = 49$ --- which is
exactly the size of the precomputed $y$-power table.

Finally
\begin{equation}
\mathrm{LINEARIZATION\_EVAL} \;=\; -A \;=\; -\nu_y(x) \pmod r .
\end{equation}

\subsubsection{The interpreter}

Most of the identity arithmetic is not emitted as Yul. It is a byte-oriented
bytecode stored inside the \emph{codehash-pinned} VK payload
($4{,}559$ bytes, $101$ instructions, plus a $178$-word constant table),
interpreted by a switch in the verifier. Because the program lives in the
pinned blob, no proof calldata can alter control flow; the guards below defend
against a \emph{mis-generated} program, not against an attacker.

The identity stream is bracketed on both sides of the program: identities at
global positions $0..3$ are emitted as direct inline Yul \emph{before} the
interpreter starts, and the trash identity (position $48$) runs as a generated
suffix \emph{after} the termination checks; the interpreted program covers
positions $4..47$.

The VM state is a program counter $\mathit{pc}$, a spilled stack pointer
$\mathit{sp}$, a cached top $\mathit{top}$ with a liveness flag
$\mathit{has\_top}$, and the two accumulators above. The instruction set
rendered into this artifact is exactly:

\begin{center}
\begin{tabular}{@{}llp{8.0cm}@{}}
\toprule
Op & Name & Effect \\
\midrule
\code{0x05} & \opc{push\_mem\_u16} & push $\code{mem}[\mathit{ptr}]$ \\
\code{0x06} & \opc{add}            & $\mathit{top} \leftarrow \mathrm{pop} + \mathit{top}$ \\
\code{0x08} & \opc{neg}            & $\mathit{top} \leftarrow -\mathit{top}$ \\
\code{0x0b} & \opc{fold\_selector} & bucket update as above \\
\code{0x0d} & \opc{mul\_const\_u8} & $\mathit{top} \leftarrow \mathit{top}\cdot C[j]$ \\
\code{0x10} & \opc{add\_mem\_u16}  & $\mathit{top} \leftarrow \mathit{top} + \code{mem}[\mathit{ptr}]$ \\
\code{0x11} & \opc{mul\_mem\_u16}  & $\mathit{top} \leftarrow \mathit{top}\cdot \code{mem}[\mathit{ptr}]$ \\
\code{0x19} & \opc{native\_permutation} & run the generated permutation block \\
\code{0x1f} & \opc{native\_lookup}      & run the generated LogUp block \\
\code{0x1b} & \opc{native\_identity}    & dispatch one of four heavy gates; each is linearised, its tail folding into $B_3$, $B_4$ or $B_9$ \\
\code{0x21} & \opc{modarith7}           & fused affine 7-limb foreign-field term \\
\bottomrule
\end{tabular}
\end{center}

\noindent
\opc{modarith7} dominates the program by weight: $4{,}331$ of the $4{,}559$
bytes, but only $16$ of the $101$ instructions. It evaluates
\begin{equation}
\mathit{acc} \;=\;
\mathrm{cond}^{[f_0]}\!\left(
c^{[f_1]}
+ \sum_{b}\sum_{i=0}^{6} c_{b,i}\,m_{b,i}
+ \sum_{b} m^{\mathrm{lhs}}_{b}\sum_{i=0}^{6} c_{b,i}\,m^{\mathrm{rhs}}_{b,i}
+ \sum_{k} c_k m_k
+ \sum_{k} c_k m^{l}_{k} m^{r}_{k}
\right),
\end{equation}
where every $c$ is a constant-table index and every $m$ a memory pointer. Its
optional \code{BILIN7\_PAIRWISE} convolution block is never instantiated in
this program (block census over the $16$ instructions:
$\mathrm{lin}=32$, $\mathrm{row}=78$, $\mathrm{pairwise}=0$,
$\mathrm{mem}=96$, $\mathrm{product}=285$), so the two tighter window guards
that protect it are dead code for this fixture. The full per-opcode census of
the decoded program is
$\opc{push\_mem\_u16}\times 12$, $\opc{add}\times 7$, $\opc{neg}\times 5$,
$\opc{fold\_selector}\times 21$, $\opc{mul\_const\_u8}\times 20$,
$\opc{add\_mem\_u16}\times 2$, $\opc{mul\_mem\_u16}\times 12$,
$\opc{native\_permutation}\times 1$, $\opc{native\_lookup}\times 1$,
$\opc{native\_identity}\times 4$, $\opc{modarith7}\times 16$ ---
$101$ in total.

\subsubsection{Structural guards}

The interpreter fails closed with \err{QuotientProgramInvalid()} on:

\begin{enumerate}[label=\textbf{Q\arabic*.},leftmargin=1.5cm]
\item \textbf{Operand window.} Every decoded memory pointer must satisfy
      $\mathit{ptr} - \mathtt{0x3680} \le_u \mathtt{0x6aa0}$, i.e.\
      $\mathit{ptr} \in [\mathtt{0x3680},\mathtt{0xa120}]$ --- one contiguous
      window from the VK payload base through the decoded-evaluation frame,
      spanning the challenge and scratch slots that sit in between. Twelve sites; two of them use the
      tighter bound $\mathtt{0x69e0}$ to leave headroom for a $7$-limb stride.
\item \textbf{Spill capacity.} $\mathit{sp} < \mathtt{0xbfa0}$ before any push
      ($54$ words).
\item \textbf{Spill underflow.} \opc{add} requires $\mathit{sp} > $ base.
\item \textbf{Selector operands.} $j < 10$ and $g \le 48$. Both bounds are
      exactly tight: bucket $10$ would land on $A$'s word, and gap $49$ would
      index past the $y$-power table.
\item \textbf{Fold liveness.} \opc{fold\_selector} requires
      $\mathit{has\_top}$.
\item \textbf{Identity boundaries.} Each native callback requires an empty
      spilled stack.
\item \textbf{Unknown opcode.} Two fail-closed \code{default} arms (the main
      switch and the heavy-gate sub-switch).
\item \textbf{Termination.} At the end, $\mathit{pc}$ must equal the program
      end exactly, $\mathit{has\_top}$ must be clear, and the spilled stack
      must be empty.
\end{enumerate}

\begin{remark}[The one asymmetry]
Constant-table \emph{indices} are the only program-derived operand with no
runtime bound. A byte index reaches $256$ slots against a $178$-slot table, so
an over-range index reads from the packed program-byte region
$[\mathtt{0x50a0},\mathtt{0x5a40}]$ --- still inside the
codehash-pinned VK payload $[\mathtt{0x3680},\mathtt{0x7900})$, never proof
calldata. It is therefore not exploitable, only silently wrong.

A second asymmetry is sharper: the three native-callback arms clear
$\mathit{has\_top}$ \emph{without checking it}, while requiring the spilled
stack to be empty. A mis-generated program that left a live cached top before a
callback would have that value silently dropped from $\nu_y(x)$, and the
end-of-program balance checks (Q8) would still pass. The rest of the VM is
carefully fail-closed about exactly this class of bug.
\end{remark}

\subsubsection{The identities themselves}

\paragraph{Permutation ($18$ columns, $6$ sets, chunk length $3$).}
With $\delta = g^{2^{32}}$ the standard Halo2 coset separator
($g=7$ the multiplicative generator of $\Fr$, $2^{32}$ the 2-adicity of
$r-1$):
\begin{align}
&L_0(x)\big(1 - Z_0(x)\big), \\
&L_{\mathrm{last}}(x)\big(Z_5(x)^{2} - Z_5(x)\big), \\
&L_0(x)\big(Z_i(x) - Z_{i-1}(\omega^{-\ell}x)\big), \quad i = 1..5, \\
&\big(1 - L_{\mathrm{last}}(x) - L_{\mathrm{blind}}(x)\big)
 \Big[
   Z_s(\omega x)\!\!\prod_{j \in \mathrm{chunk}(s)}\!\!\big(v_j + \beta\sigma_j + \gamma\big)
   -
   Z_s(x)\!\!\prod_{j \in \mathrm{chunk}(s)}\!\!\big(v_j + \beta\,\delta^{\,j} x + \gamma\big)
 \Big],
\end{align}
for $s = 0..5$. The column values $v_j$ are heterogeneous: seventeen are
advice or fixed evaluations, and one is $\mathrm{inst}(x)$ from Phase 4 --- an
instance column participating in the permutation argument.

\paragraph{Lookup (LogUp, $2$ arguments, input widths $4$ and $1$).}
Writing $f_i = \theta\text{-compression of input }i$, $t$ for the compressed
table expression, $m$ for the multiplicity evaluation, $h$ for the helper
evaluation and $Z^{L}$ for the accumulator:
\begin{align}
&\big(L_0(x) + L_{\mathrm{last}}(x)\big)\, Z^{L}(x), \\
&h(x)\prod_{i=0}^{3}\big(f_i + \beta\big) \;-\; \sum_{i=0}^{3}\prod_{k \ne i}\big(f_k + \beta\big)
&&\Big(\text{i.e. } h = \sum_i \tfrac{1}{f_i+\beta}\Big), \\
&\big(1 - L_{\mathrm{last}}(x) - L_{\mathrm{blind}}(x)\big)
 \Big[\big(Z^{L}(\omega x) - Z^{L}(x) - s\!\cdot\!\textstyle\sum_c h_c(x)\big)\big(t + \beta\big) + m(x)\Big].
\end{align}
The middle identity is evaluated with a prefix/suffix product pass so the
$\sum_i \prod_{k\ne i}$ costs $O(w)$ rather than $O(w^{2})$ multiplications.
The width-$4$ formulas describe argument $0$; argument $1$ has a single
$\theta$-compressed input, its helper identity degenerating to
$h(x)\,(f+\beta) - 1$ with no product pass.

\paragraph{Trash.}
Eight rows are batched by $\tau$ into $C = \sum_{k=0}^{7} R_k\,\tau^{\,7-k}$
and folded as $C - (1-\mathrm{sel})\cdot \mathrm{trash}(x)$, where
$\mathrm{sel}$ is a fixed selector evaluation and $\mathrm{trash}(x)$ the
trashcan commitment's opening. Because $\tau$ was squeezed \emph{before} the
trashcan commitment was absorbed, the prover cannot choose the commitment to
suit the batch.

% ---------------------------------------------------------------------------
\subsection{Phase 6 --- Linearisation scalars}
% ---------------------------------------------------------------------------

Rather than materialise $\mathcal{L}$ with its own MSM, the verifier computes
two scalars and lets the final MSM expand \eqref{eq:lincom} in place:
\begin{equation}
x_{\mathrm{split}} = x^{\,n-1},
\qquad
1 - x^{n}.
\end{equation}
Both come from one $k$-step walk maintaining the invariant
$\big(x^{2^{i}},\,x^{2^{i}-1}\big)$:
\[
x^{2^{i+1}-1} = \big(x^{2^{i}-1}\big)^{2}\cdot x,
\qquad
x^{2^{i+1}} = \big(x^{2^{i}}\big)^{2}.
\]
The two scalars are parked in the first two words of the slot named
\mptr{QUOTIENT\_MPTR}, which on this path is scalar metadata rather than the
$\Gone$ commitment its name and declaration comment suggest.

% ---------------------------------------------------------------------------
\subsection{Phase 7 --- The KZG multi-open}
\label{sec:pcs}
% ---------------------------------------------------------------------------

\subsubsection{Rotation points and point sets}

Four distinct rotations occur in this circuit, $\{-\ell,\,-1,\,0,\,+1\}$ with
$\ell = 10$. Their evaluation points are
\begin{equation}
z_0 = x\,\omega^{-10},\qquad
z_1 = x\,\omega^{-1},\qquad
z_2 = x,\qquad
z_3 = x\,\omega .
\end{equation}
Queries are grouped into five point sets, each the set of rotations at which
one group of commitments is opened:
\begin{equation}
S_0=\{z_2\},\quad
S_1=\{z_2,z_1\},\quad
S_2=\{z_2,z_3\},\quad
S_3=\{z_2,z_3,z_1\},\quad
S_4=\{z_2,z_3,z_0\}.
\end{equation}

\subsubsection{The \texorpdfstring{$x_1$}{x1} batching within a set}

For each set $s$ and each rotation $t \in S_s$, the claimed evaluations of the
commitments in that set are batched by powers of $x_1$:
\begin{equation}
\rho_{s,t} \;=\; \sum_{j} \trunc\!\left(x_1^{\,j}\right)\, e_{s,j,t},
\label{eq:rho}
\end{equation}
where $j$ indexes the commitments of set $s$ in generated order and $e_{s,j,t}$
is that commitment's alleged evaluation at $t$. The stored power table holds
$43$ entries: $\trunc(x_1^{i})$ for $1 \le i \le 42$, and an untruncated
literal $1$ at index $0$.

Set $S_0$ is the large one --- $43$ evaluation terms over $42$ commitments,
because one enumerated commitment is the $\Gone$ identity and is skipped on the
commitment side while its evaluation slot is kept. The $43$rd evaluation term
is $-\nu_y(x)$: the linearisation claim \eqref{eq:linclaim} enters the
multi-open as an ordinary rotation-$0$ query.

\subsubsection{Per-set quotient evaluation}

Let $Z_s(X) = \prod_{t \in S_s}(X-t)$ and let $\Lambda_s$ be the interpolating
polynomial through $\{(t,\rho_{s,t})\}_{t \in S_s}$. The prover supplies
$q_s$, the claimed evaluation at $x_3$ of the per-set quotient. The verifier
recomputes
\begin{equation}
\mathrm{ev}_s
\;=\;
\frac{q_s - \Lambda_s(x_3)}{Z_s(x_3)}
\;=\;
\frac{q_s}{\prod_{t\in S_s}(x_3-t)}
\;-\;
\sum_{t \in S_s}
\frac{\rho_{s,t}}{(x_3-t)\prod_{u \in S_s\setminus\{t\}}(t-u)},
\label{eq:evs}
\end{equation}
the right-hand form being the barycentric expansion the code actually emits.
For $|S_s| = 1$ this degenerates to $(q_0-\rho_{0})/(x_3-x)$.

Each set's differences and Lagrange denominators are inverted with a
\emph{single} \code{modexp}: a running product $\mathrm{bp}$ is built over all
$2|S_s|$ factors, inverted once by Fermat, and unwound (for $|S_s| = 1$ the
single difference is inverted directly; there is no product to unwind). The helper rejects a
zero or non-canonical input with \err{NonCanonicalScalar()}, so a degenerate
$x_3$ fails closed rather than producing zero inverses.

\subsubsection{The \texorpdfstring{$x_2$}{x2} fold}

\begin{equation}
f_{\mathrm{eval}} \;=\; \sum_{s=0}^{4} \mathrm{ev}_s\; x_2^{\,s},
\end{equation}
emitted as a Horner fold in decreasing $s$. $x_2$ is used at full precision.

\subsubsection{The \texorpdfstring{$x_4$}{x4} batch and the fused final MSM}

\begin{align}
v &\;=\; \sum_{s=0}^{4} q_s\,\trunc\!\left(x_4^{\,s}\right)
        \;+\; f_{\mathrm{eval}}\,\trunc\!\left(x_4^{\,5}\right),
\label{eq:v}\\[1mm]
\mathrm{final\_com}
&\;=\;
\sum_{s=0}^{4} \trunc\!\left(x_4^{\,s}\right)
   \sum_{j} \trunc\!\left(x_1^{\,j}\right) C_{s,j}
\;+\;
\trunc\!\left(x_4^{\,5}\right) F,
\label{eq:finalcom}
\end{align}
where $C_{s,j}$ ranges over the commitments of set $s$, \emph{except} that the
linearised commitment in $S_0$ is not a single point: it is expanded in place
via \eqref{eq:lincom} into $4$ quotient-limb terms with scalars
$\trunc(x_1^{\,j_\star})\,(1-x^{n})\,x_{\mathrm{split}}^{\lambda}$ and $10$
selector terms with scalars $\trunc(x_1^{\,j_\star})\,B_j$.

The arithmetic works out to exactly $78$ terms:
\[
\underbrace{41}_{S_0,\ \text{ordinary}}
+ \underbrace{4}_{Q_\lambda} + \underbrace{10}_{S_j}
+ \underbrace{3+3+11+5}_{S_1..S_4}
+ \underbrace{1}_{F}
\;=\; 78,
\]
staged as $78$ consecutive $(\text{point},\text{scalar})$ pairs of
$\mathtt{0xa0}$ bytes and consumed by a single \code{G1MSM} call at the exact
scheduled cost for $78$ pairs.

\begin{remark}[Truncated batching challenges]
$x_1$, $x_4$ and $x_3$ enter the batch truncated to $128$ bits. This is the
\code{midnight-proofs} truncated-challenge convention: it raises the soundness
error of the batching step from ${\approx}\,\deg/r$ to
${\approx}\,\deg/2^{128}$, which is still negligible --- but a reviewer meeting
\code{and(acc, 0xffff...ffff)} on a challenge-derived scalar, with no comment in
scope, is meeting the exact pattern that usually signals a soundness bug.
\end{remark}

% ---------------------------------------------------------------------------
\subsection{Phase 8 --- Pairing, and the randomised accumulator batch}
\label{sec:pairing}
% ---------------------------------------------------------------------------

\subsubsection{The KZG equation}

For a single opening of $\mathrm{final\_com}$ at $x_3$ with claimed value $v$
and proof $\pi$, KZG requires
\begin{equation}
\pair\big(\mathrm{final\_com} - v\,G + x_3\,\pi,\ [1]_2\big)
\;=\;
\pair\big(\pi,\ [s]_2\big).
\label{eq:kzg}
\end{equation}
The verifier stages
\begin{equation}
R \;=\; \mathrm{final\_com} - v\,G + x_3\,\pi,
\qquad
L \;=\; \pi,
\end{equation}
building $R$ with one \code{G1MSM} for $(-v)G$, one \code{G1ADD}, one
\code{G1MSM} for $x_3\pi$, and one \code{G1ADD}. Every staging call --- these,
the $78$-pair MSM before them, and the accumulator-batch calls below --- folds
into one success flag that is surfaced as \err{PrecompileFailed()}
\emph{before} the pairing helper runs.

\begin{remark}[Naming]
$R$ and $L$ are stored in slots named \mptr{PAIRING\_RHS\_MPTR} and
\mptr{PAIRING\_LHS\_MPTR}, but those names follow the historical dual-MSM
accumulator (left $=\pi$, right $=$ combined), \emph{not} the pairing argument
order. They are passed to the pairing helper swapped, and the source needs a
thirteen-line comment to say so.
\end{remark}

\subsubsection{Batching in the IVC accumulator}

The public accumulator asserts its own pairing equation over
$(\mathrm{acc}_L,\mathrm{acc}_R)$. Multiplying the two equations together
would be unsound --- two wrong equations can cancel --- so the verifier draws a
local randomiser \emph{after} all four $\Gone$ points are fixed:
\begin{equation}
\alpha \;=\;
\mathrm{keccak256}\big(
D \,\|\, \mathrm{vk\_digest} \,\|\, R \,\|\, L \,\|\, \mathrm{acc}_R \,\|\, \mathrm{acc}_L
\big) \bmod r,
\qquad
\alpha \leftarrow 1 \text{ if } \alpha = 0,
\label{eq:alpha}
\end{equation}
over a $576$-byte preimage, where
$D = \mathtt{0x00}^{3}\,\|\,\texttt{"pairing-batch-acc-kzg"}\,\|\,\mathtt{0x00}^{8}$.
Including $\mathrm{vk\_digest}$ binds $\alpha$ to the verifying key directly
rather than transitively through the points.

Then
\begin{equation}
R \leftarrow R + \alpha\,\mathrm{acc}_R,
\qquad
L \leftarrow L + \alpha\,\mathrm{acc}_L,
\end{equation}
and if either original equation is false, the combined equation holds for at
most one $\alpha \in \Fr$. The $\alpha = 0$ substitution matters: without it a
zero draw would silently delete the accumulator equation.

\subsubsection{The accepting gate}

The pairing helper stages
$[\,R' \mid [1]_2 \mid L' \mid [-s]_2\,]$ --- $\mathtt{0x300}$ bytes, the two
$\Gtwo$ bases read from the pinned VK payload --- and calls \code{0x0f} at the
exact two-pair schedule cost. It accepts iff the returned word is \emph{equal
to} $1$:
\begin{equation}
\pair\big(R',\,[1]_2\big)\cdot \pair\big(L',\,[-s]_2\big) \;=\; 1_{\GT}
\quad\Longleftrightarrow\quad
\pair\big(R',\,[1]_2\big) = \pair\big(L',\,[s]_2\big).
\end{equation}
A failed \code{staticcall} or a short return is \err{PrecompileFailed()}; a
returned $0$ is \err{ProofRejected()}. On acceptance the function writes $1$ to
memory and returns from assembly; there is no path that returns without a
completed pairing check.

% ===========================================================================
\section{Complete check catalogue}
\label{sec:catalogue}
% ===========================================================================

Every runtime rejection in the two artifacts, by phase. Conditions are stated
as the \emph{accept} condition. All eight selectors are
\code{bytes4(keccak256("Name()"))} of the declared errors.

\begin{center}
\begin{tabular}{@{}ll@{}}
\toprule
Error & Selector \\
\midrule
\err{BadCalldataShape()}       & \code{0x1b99e37c} \\
\err{VkMismatch()}             & \code{0xa447d73e} \\
\err{NonCanonicalScalar()}     & \code{0x77530042} \\
\err{BadPointEncoding()}       & \code{0xf27905ec} \\
\err{PrecompileFailed()}       & \code{0x84e81692} \\
\err{ProofRejected()}          & \code{0xc3b0d8cd} \\
\err{QuotientProgramInvalid()} & \code{0x3cc81b89} \\
\err{MemoryLayoutViolated()}   & \code{0xc9888d23} \\
\bottomrule
\end{tabular}
\end{center}

\renewcommand{\arraystretch}{1.2}
\begin{longtable}{@{}p{0.8cm}p{7.1cm}p{6.6cm}@{}}
\caption{Runtime checks, in execution order.}\label{tab:checks}\\
\toprule
\# & Accept condition & On failure \\
\midrule
\endfirsthead
\multicolumn{3}{@{}l}{\small\itshape Table \ref{tab:checks}, continued}\\
\toprule
\# & Accept condition & On failure \\
\midrule
\endhead
\bottomrule
\endfoot

\multicolumn{3}{@{}l}{\textbf{Deployment (constructor)}}\\
D1 & $\code{mload}(\mathtt{0x40}) \le \mathtt{0x1000}$ in the creation frame & \err{MemoryLayoutViolated()} \\
D2 & \opc{mcopy} round-trip returns the stored word & bare revert \\
D3 & \code{modexp} succeeds, returns $32$ bytes, and $2\cdot\mathrm{res}\equiv 1 \bmod r$ & bare revert \\
D4 & $\code{G1ADD}(\Ident,\Ident)=\Ident$ and $\code{G1ADD}(G,G)=[2]G$ & bare revert \\
D5 & $\code{G1MSM}([2]\cdot G)=[2]G$ & bare revert \\
D6 & \code{G1MSM} \emph{rejects} an on-curve, off-subgroup point & bare revert \\
D7 & $\code{PAIRING}(G,H;-G,H)$ returns exactly $1$ & bare revert \\
D8 & $\code{PAIRING}(G,H;G,H)$ returns exactly $0$ & bare revert \\
D9 & $78$-pair all-identity \code{G1MSM} at the runtime gas bound returns $128$ bytes encoding $\Ident$ & bare revert \\
D10 & all-identity two-pair \code{PAIRING\_CHECK} returns exactly $1$ & bare revert \\
D11 & $\code{code.length}=17{,}025$ and $\code{codehash}=\mathtt{0xe68d}\ldots$ & \code{require("invalid vk")} \\
\addlinespace
\multicolumn{3}{@{}l}{\textbf{Phase 0 --- entry}}\\
0.1 & $\code{cd}[\mathtt{0x04}]=\mathtt{0x40}$ and $\code{cd}[\mathtt{0x24}]=\mathtt{0x1ec0}$ & \err{BadCalldataShape()} \\
0.2 & $\code{mload}(\mathtt{0x40}) \le \mathtt{0x1000}$ & \err{MemoryLayoutViolated()} \\
\addlinespace
\multicolumn{3}{@{}l}{\textbf{Phase 1 --- VK and envelope}}\\
1.1 & $\code{extcodesize}(vk)=17{,}025$ & \err{VkMismatch()} \\
1.2 & $\code{extcodehash}(vk)=\mathtt{0xe68d}\ldots$ & \err{VkMismatch()} \\
1.3 & header words equal $19,\,20,\,1,\,11,\,7,\,56$ & \err{VkMismatch()} \\
1.4 & $\code{proof.length}=\mathtt{0x1e60}$ & \err{BadCalldataShape()} \\
1.5 & $\code{instances.length}=19$ & \err{BadCalldataShape()} \\
1.6 & $\code{calldatasize}=\mathtt{0x2144}$ & \err{BadCalldataShape()} \\
\addlinespace
\multicolumn{3}{@{}l}{\textbf{Phase 2 --- public accumulator (2.1--2.3 per coordinate $\times 4$; 2.4--2.6 per point $\times 2$)}}\\
2.1 & $w_0 < 2^{224}$ and $w_1 < 2^{168}$ (no unused high bits) & \err{BadPointEncoding()} \\
2.2 & $V \le p-1$ in split hi/lo form & \err{BadPointEncoding()} \\
2.3 & $c_{\mathrm{hi}} < 2^{128}$ & \err{BadPointEncoding()} \\
2.4 & the identity is encoded only by the exact four-word sentinel & \err{BadPointEncoding()} \\
2.5 & a non-identity point does not decode to $(0,0)$ & \err{BadPointEncoding()} \\
2.6 & if $x$ carries the identity flag, $y$ decodes to $0$ & \err{BadPointEncoding()} \\
2.7 & $\code{G1MSM}([1]\cdot \mathrm{acc}_L)$ succeeds, returns $\mathtt{0x80}$ bytes & \err{PrecompileFailed()} \\
2.8 & $\code{G1MSM}([1]\cdot \mathrm{acc}_R)$ succeeds, returns $\mathtt{0x80}$ bytes & \err{PrecompileFailed()} \\
\addlinespace
\multicolumn{3}{@{}l}{\textbf{Phase 3 --- transcript}}\\
3.1 & every instance word $< r$ & \err{NonCanonicalScalar()} \\
3.2 & every proof evaluation $< r$ ($102$ of them) & \err{NonCanonicalScalar()} \\
3.3 & every $q_s$ $< r$ ($5$ of them) & \err{NonCanonicalScalar()} \\
3.4 & every absorbed $\Gone$: $x_{\mathrm{hi}} < 2^{128}$, $y_{\mathrm{hi}} < 2^{128}$ & \err{BadPointEncoding()} \\
3.5 & every absorbed $\Gone$: $x < p$ and $y < p$ & \err{BadPointEncoding()} \\
3.6 & the parser cursor lands exactly on the instances length word & \err{BadCalldataShape()} \\
\addlinespace
\multicolumn{3}{@{}l}{\textbf{Phase 4 --- domain evaluation}}\\
4.1 & every batched denominator $< r$ & \err{ProofRejected()} \\
4.2 & the product of all $30$ denominators is nonzero, i.e.\ $x$ is not a domain point and $x^{n}\neq 1$ & \err{ProofRejected()} \\
4.3 & \code{modexp} succeeds and returns $32$ bytes & \err{PrecompileFailed()} \\
\addlinespace
\multicolumn{3}{@{}l}{\textbf{Phase 5 --- quotient VM}}\\
5.1 & every decoded pointer lies in $[\mathtt{0x3680},\mathtt{0xa120}]$ (tighter, $\mathtt{0xa060}$, for stride bases) & \err{QuotientProgramInvalid()} \\
5.2 & spill pointer $< \mathtt{0xbfa0}$ before a push & \err{QuotientProgramInvalid()} \\
5.3 & \opc{add} has a spilled operand & \err{QuotientProgramInvalid()} \\
5.4 & selector index $< 10$ & \err{QuotientProgramInvalid()} \\
5.5 & selector gap $\le 48$ & \err{QuotientProgramInvalid()} \\
5.6 & \opc{fold\_selector} has a live cached top & \err{QuotientProgramInvalid()} \\
5.7 & each native callback is entered with an empty spilled stack & \err{QuotientProgramInvalid()} \\
5.8 & opcode is one of the eleven rendered; callback index is one of four & \err{QuotientProgramInvalid()} \\
5.9 & at end: $pc$ equals the program end, no live top, empty spilled stack & \err{QuotientProgramInvalid()} \\
\addlinespace
\multicolumn{3}{@{}l}{\textbf{Phase 7 --- multi-open}}\\
7.1 & each per-set inversion input is nonzero and $< r$ & \err{NonCanonicalScalar()} \\
7.2 & each \code{modexp} succeeds and returns $32$ bytes & \err{PrecompileFailed()} \\
7.3 & the $78$-pair \code{G1MSM} succeeds and returns $\mathtt{0x80}$ bytes & \err{PrecompileFailed()} \\
\addlinespace
\multicolumn{3}{@{}l}{\textbf{Phase 8 --- pairing}}\\
8.1 & every staging \code{G1MSM}/\code{G1ADD} succeeds with the right return size & \err{PrecompileFailed()} \\
8.2 & the pairing \code{staticcall} succeeds and returns $32$ bytes & \err{PrecompileFailed()} \\
8.3 & the returned word equals $1$ & \err{ProofRejected()} \\
\end{longtable}

\noindent
Three further rejection sites are rendered but unreachable for this fixture,
listed here so the table can claim completeness: the \code{ec\_pairing} entry
guard (\err{ProofRejected()} on a cleared success flag, pre-empted by the
\err{PrecompileFailed()} conversion immediately before the call), the
redundant post-pairing success guard (the source itself notes the redundancy),
and \code{batch\_invert}'s reversed-range and single-element paths, which its
only call site --- the fixed $30$-word Lagrange run --- can never take.

% ===========================================================================
\section{What the contract does not check}
\label{sec:notchecked}
% ===========================================================================

The list below is not a list of defects. Each item is a deliberate delegation,
and knowing where the obligation actually lives is the difference between an
auditable design and a gap.

\begin{enumerate}[label=\textbf{N\arabic*.},leftmargin=1.4cm]

\item \textbf{Curve and subgroup membership of proof commitments.}
\code{common\_uncompressed\_g1} checks padding and $\Fp$ range only. Membership
is enforced by EIP-2537 itself: the generator's \code{ProtocolPlan::validate}
rejects any plan in which an absorbed commitment would not later reach a
\code{G1MSM} or the pairing, and those precompiles perform the check. The
constructor's negative \code{G1MSM} probe (D6) is what makes that delegation
safe on a given chain. Note the consequence for error attribution: an
off-curve accumulator point surfaces as \err{PrecompileFailed()}, not
\err{BadPointEncoding()}.

\item \textbf{The identity as a commitment.} Nothing rejects
$\Ident$ for an absorbed proof commitment. It is a valid EIP-2537 input and a
valid group element; a prover gains nothing by sending it.

\item \textbf{Domain constants.} $n^{-1}$, $\omega$, $\omega^{-1}$ and
$\omega^{-\ell}$ are used without any algebraic cross-check
($\omega\cdot\omega^{-1}=1$, $\mathrm{ord}(\omega)=2^{k}$). Their correctness
is covered by the VK codehash pin alone.

\item \textbf{The trusted setup.} $[-s]_2$ is unattested data. Nothing on chain
relates it to any ceremony transcript; the \code{BUILD\_ID} preimage is the
only place an SRS fingerprint appears, and it is published off chain.

\item \textbf{The $\mathrm{vk\_digest}$ itself.} It is absorbed first into the
transcript, but never checked against the verifying key it labels. It is a
binding tag, not a commitment the contract verifies.

\item \textbf{Semantics of the public instances.} \code{verifyProof} checks
only that the proof verifies under the pinned key. Binding those $19$ words to
state roots, program identifiers, chain/domain separation or anti-replay is an
application-wrapper obligation.

\item \textbf{Word alignment of the batch-inversion range}, and the capacity of
its scratch window. Both are generation-time invariants. The scratch abuts the
denominator run with exactly zero slack, and an overrun would be silent: the
\code{modexp} would still succeed and wrong inverses would be returned without
a revert.

\item \textbf{Fold count in the quotient program.} The epilogue validates the
program counter and the stack, but nothing counts identities. The tail
exponents encode $N=49$; a bytecode whose fold count disagreed would misalign
every selector bucket, and no runtime check would notice.

\item \textbf{Constant-table index range} in the quotient VM (see the remark in
Section~\ref{sec:qvm}). Out of range reads pinned VK bytes, never calldata.

\end{enumerate}

% ===========================================================================
\section{Notes for a reader of the source}
\label{sec:reading}
% ===========================================================================

Three conventions are worth internalising before opening the \code{.sol}.

\paragraph{Absolute addressing, phase-reused.}
The verifier does not use Solidity's free-memory pointer. It allocates from one
arena of absolute addresses, and several regions deliberately share a base
because their lifetimes are disjoint. Six regions alias $\mathtt{0x1000}$
alone. The reuse is checked at \emph{generation} time by an allocator that
rejects two live regions sharing an address; nothing at run time re-derives it.
The practical consequence: seeing one address on several constants means
``reused in sequence,'' never ``written twice at once'' --- but a reader
verifying that locally cannot, and must trust the generator.

The sharpest instance is $\mathtt{0xb280}$, which is
(i) the declared constant \mptr{PCS\_FINAL\_MSM\_MPTR},
(ii) the $y$-Horner accumulator $A$ throughout Phase 5,
(iii) the base of three staged eval-source tables in Phase 7, and
(iv) the $\mathtt{0x30c0}$-byte final MSM staging buffer.
It is written symbolically in (i) and (iv) and as a bare literal in (ii) and
(iii).

\paragraph{Deferred failure, then a section boundary.}
Many checks fold into a \code{success} flag rather than reverting on the spot,
and a single guard at the end of the block converts it to a typed revert. Where
a chain fault and a proof fault could both clear the same flag, a second
\code{precompile\_failed} flag disambiguates them --- this is the ``MF-4''
pattern that recurs in the batch inversion, the accumulator validator and the
Lagrange block.

\paragraph{Guards for the generator, not the prover.}
A large fraction of the runtime guards --- the entire quotient-VM guard set,
the range checks in the batch inversion, the parser-closure check --- cannot be
triggered by any calldata, because the thing they validate is
codehash-pinned. They defend against a mis-generated artifact. The source does
not distinguish these from attacker-reachable guards, so a reviewer must
classify each one by hand.

% ===========================================================================
\section{Provenance of this document}
% ===========================================================================

Everything above was read from
\code{fixtures/moonlight-wrap/Halo2Verifier.sol} (render provenance pinned by
the fixture README's source-commit table) and
\code{fixtures/moonlight-wrap/Halo2VerifyingKey.sol}. Structural counts
(instruction census, block census, point-set membership, MSM width, tail
exponents) were obtained by decoding the packed quotient program in the VK
payload against the interpreter in the verifier, not by consulting the Rust
generator.

A companion document, \code{docs/spec/Halo2VerifierSpec.tex}, covers the same
two artifacts line by line with quoted excerpts and a numbered findings
register; \code{docs/spec/check\_excerpts.py} re-verifies its citations against
the fixtures. This document is deliberately independent of it: where the two
agree, the agreement is a cross-check.

Line numbers are omitted here on purpose. Re-rendering the verifier invalidates
them, whereas the structure and the algebra above survive.

\end{document}
